\documentclass[a4paper]{article}
\usepackage[utf8]{inputenc}
\usepackage[T2A]{fontenc}
\usepackage[english,russian]{babel}
\usepackage[12pt]{extsizes}
\usepackage{hyperref}

\usepackage[left=4cm,right=3cm,top=3cm,bottom=3cm,bindingoffset=0cm]{geometry}
\tolerance=10000 % разреженность строки

% \usepackage{indentfirst} % красная строка для первой строки абзаца
\setlength{\parindent}{0.9cm} % длина отступа красной строки
\setlength{\parskip}{0.3em} % интервал между абзацами
\linespread{1.1} % интерлиньяж
\renewcommand{\labelitemi}{\textendash} % дефис как маркер пунктов списка

\title{Выпускная квалификационная работа}
\author{Амир Муллагалиев}
\date{Май 2021}

\begin{document}

\begin{titlepage}
    \begin{center}
        Филиал Московского Государственного Университета\\
        им. М.В.Ломоносова в городе Ташкенте
        \vspace{0.4cm}

        Факультет прикладной математики и информатики\\
        Кафедра прикладной математики и информатики
        \vfill


        Муллагалиев Амир Зиннурович
        \vfill

        {\Large\textsc{Выпускная квалификационная работа}}\\
        на тему: «Безопасность и оптимизация процесса проверки работ в системе дистанционного обучения»
        \bigskip

        по направлению 01.03.02 «Прикладная математика и информатика»
    \end{center}
\vfill

\begin{flushleft}
Выпускная квалификационная работа рассмотрена и рекомендована к защите\\
Руководитель филиала\\
доцент Строгалов А.С.\\
Научный руководитель\\
к.ф.м.н, в.н.с. Алисейчик Павел Александрович
\end{flushleft}
\vspace{0.2cm}
\begin{flushright}
«\underline{\hspace{0.7cm}}» \underline{\hspace{3cm}} 2021 г.
\end{flushright}
\vspace{0.7cm}

\begin{center}
Ташкент 2021 г.
\end{center}
\end{titlepage}



\newpage
\begin{abstract}
В данной работе описаны результаты, достигнутые в ходе расширения функциональности системы дистанционного обучения «МГУ Контест», усовершенствования процесса проверки работ студентов и обеспечения безопасности при компиляции и запуске пользовательских программ в системе.

Внедрено решение для редактирования пользовательского текста. Произведено перепроектирование существующей программной структуры системы для реализации в системе новых видов выдаваемых для решения студентами задач. Внесены структурные решения и описан пользовательский интерфейс, реализующие новые виды выдаваемых задач и способы взаимодействия с ними пользователей. Применён инструмент для изоляции процесса компиляции и запуска пользовательских программ в системе с целью обеспечения безопасности на сервере.
\end{abstract}



\newpage
\tableofcontents



\newpage
\section{Введение}



\newpage
\section{Форма редактирования пользовательского текста}



\newpage
\section{Реализация новых видов выдаваемых для решения студентами задач}



\newpage
\subsection{Перепроектирование программной структуры системы}



\newpage
\subsection{Структурные решения и пользовательский интерфейс, реализующие новые виды выдаваемых задач}



\newpage
\section{Изоляция процесса компиляции и запуска пользовательских программ}



\newpage
\section{Заключение}
В рамках данной выпускной квалификационной работы внедрено решение для редактирования пользовательского текста, произведено перепроектирование существующей программной структуры системы для реализации в системе новых видов выдаваемых для решения студентами задач, внесены структурные решения и описан пользовательский интерфейс, реализующие новые виды выдаваемых задач и способы взаимодействия с ними пользователей, применён инструмент для изоляции процесса компиляции и запуска пользовательских программ в системе с целью обеспечения безопасности на сервере.



\newpage
\section{Список использованных источников и литературы}
\begingroup
\renewcommand{\section}[2]{}
\begin{thebibliography}{3}
    \bibitem{intsys}
    Алисейчик П.А., Вашик К., Ж.Кнап, Кудрявцев В.Б., С.Г. Шеховцов, Строгалов А.С. Моделирование процесса обучения // Интеллектуальные системы, Т.10, 2006 г. 85 стр.

    \bibitem{bekta}
    Алисейчик П.А., Строгалов А.С., Бекташев Р.А. О дистанционном образовании -- пример реализации и перспективы // Интеллектуальные системы. Теория и приложения (ранее: Интеллектуальные системы по 2014, № 2, ISSN 2075-9460), Т.20, № 3, с. 127-133

    \bibitem{django}
    Документация фреймворка Django\\
    \url{https://docs.djangoproject.com/en/3.0}

    \bibitem{python}
    Документация Python\\
    \url{https://docs.python.org/3}

    \bibitem{contestrepo}
    GitHub-репозиторий проекта «МГУ Контест»\\
    \url{https://github.com/ruslanbektashev/contest}

    \bibitem{contest}
    Сайт «МГУ Контест»\\
    \url{https://contest.msu-dev.ru/}

    \bibitem{ckeditor}
    Сайт «CKEditor»\\
    \url{https://ckeditor.com/}

    \bibitem{docker}
    Документация-справочник по инструменту «Docker»\\
    \url{https://docs.docker.com/reference/}
\end{thebibliography}
\endgroup
\newpage



\end{document}
